\section{Problem Statement and Research Objectives}\label{sec:problem-statement-and-research-objectives}

Existing work on MSAR, Digital Twins, and HCI in Supervisory Control systems has revealed critical gaps in the literature.
Unmanned assets have been shown, by~\cite{7761074}, to improve coverage and responsiveness, while also raising concerns about operator overload as fleets grow.
Many Digital Twin implementations in the Maritime domain emphasize physical modelling and asset performance, as elements like cognitive demand and supporting decision-processes are de-prioritized.
Studies in HCI have also highlighted the importance of ensuring a balance between automation and human control, as well as the need to account for operator workload and situational awareness in designing supervisory interfaces.

This project aims to address these gaps by developing a Digital Twin-inspired, layered architecture that integrates swarm simulation, path planning, and decision support tools for a drone swarm in an MSAR-situation.
The mission control interface will be designed with a priority placed on workload migration, applying design principles like the Proximity Compatibility Principle, and introducing features like structured alert hierarchies to reduce extraneous cognitive load and improve situational awareness.
Finally, it will also define a series of subjective and objective metrics that can be used in future human-subject studies to quantify the effectiveness of our system.


